\documentclass[11pt,a4paper]{article}

% --- Packages ---
\usepackage[utf8]{inputenc}
\usepackage[margin=1in]{geometry}
\usepackage{amsmath,amssymb}
\usepackage{booktabs}
\usepackage{hyperref}
\usepackage{xcolor}
\usepackage{tcolorbox}
\usepackage{enumitem}

% --- Styling ---
\hypersetup{
    colorlinks=true,
    linkcolor=teal,
    urlcolor=teal
}

\newtcolorbox{plainbox}[1]{
    colback=teal!6!white,
    colframe=teal!70!black,
    fonttitle=\bfseries,
    title=#1,
    arc=2mm
}

\newtcolorbox{alertbox}[1]{
    colback=amber!6!white,
    colframe=orange!85!black,
    fonttitle=\bfseries,
    title=#1,
    arc=2mm
}

\title{\textbf{\LARGE Clover: How It Works From The Ground Up}\\[0.4em]
\large A Simple, Plain-English Guide to Satellite Tree Counting, Canopy Measurement, and Carbon Math}
\author{\textbf{Soumodeep Karmakar} \\ Flora Carbon Submission}
\date{September 2026}

\begin{document}

\maketitle

\begin{abstract}
This guide explains every single concept, algorithm, and formula behind \textbf{Clover} in plain, intuitive English. No dense academic jargon, no unexplained equations. If you want to understand how a satellite picture of trees turns into an accurate, auditable carbon number---and what the math is actually doing behind the scenes---this document breaks it down step-by-step from scratch.
\end{abstract}

\tableofcontents
\newpage

% ==============================================================================
\section{The Big Picture: What Problem Are We Solving?}
% ==============================================================================

\subsection{Why Forest Carbon Matters}
Companies and governments around the world want to offset their emissions by planting or protecting trees. But in carbon markets, you can only sell a carbon credit if you can \textbf{prove} how much carbon is actually stored in that forest.

\subsection{The Old Way vs. The AI Trap}
\begin{itemize}
    \item \textbf{The Old Way (Manual Forestry)}: Foresters walk into the woods with measuring tapes and clipboards. They measure tree trunks one by one. It is slow, costs a fortune, and you can only inspect a tiny fraction ($1\text{--}2\%$) of the forest.
    \item \textbf{The Lazy AI Trap}: Companies take satellite photos, run generic object detection models, draw rectangles on green spots, add up the numbers, and claim 99.9\% accuracy. 
\end{itemize}

\begin{alertbox}{Why Most Satellite Carbon Tools Are Misleading}
\begin{enumerate}
    \item \textbf{Rectangles aren't trees}: Trees are round; boxes have empty corners. Rectangles overestimate tree area by over $21\%$.
    \item \textbf{Overlap double-counting}: In a real forest, tree branches overlap. Adding up individual boxes counts the shared area twice, sometimes claiming $120\%$ canopy cover!
    \item \textbf{Fictitious certainty}: Estimating carbon from satellite pictures has natural error margins ($\pm 10\text{--}15\%$). Claiming an exact number with zero error bars is scientifically dishonest.
\end{enumerate}
\end{alertbox}

\texttt{Clover} was built to solve these three flaws: it detects individual trees, removes overlaps using real geometry, applies peer-reviewed carbon scaling equations, and honestly shows you its margin of error.

% ==============================================================================
\section{Step 1: The Camera and The Pixels (GSD)}
% ==============================================================================

\subsection{What is Ground Sampling Distance (GSD)?}
When a satellite takes a photo of Earth, every pixel represents a specific patch of real ground. This is called the \textbf{Ground Sampling Distance (GSD)}.
\begin{equation}
\text{GSD} = \text{Real-world distance covered by one pixel (e.g., } 25\,\text{cm/pixel)}
\end{equation}

\begin{itemize}
    \item \textbf{Drone / UAV ($\text{GSD} < 10\,\text{cm}$)}: Super sharp. You can see individual branches and leaf clusters.
    \item \textbf{Commercial Satellites ($\text{GSD} = 25\text{--}50\,\text{cm}$, e.g., WorldView)}: Sharp enough to clearly see single tree crowns and their circular shapes.
    \item \textbf{Free Satellites ($\text{GSD} = 10\,\text{meters}$, e.g., Sentinel-2)}: \textbf{Too blurry for single trees!} A single pixel is $10\,\text{m} \times 10\,\text{m}$ (the size of an entire house). Three or four trees are mashed into one green square.
\end{itemize}

\begin{plainbox}{Our Built-in Safety Rule}
If someone uploads imagery with $\text{GSD} > 1.0\,\text{meter}$, \texttt{Clover} displays an explicit warning. At that scale, an AI cannot see individual tree crowns; it would just be hallucinating numbers.
\end{plainbox}

% ==============================================================================
\section{Step 2: World Maps and Metric Coordinates}
% ==============================================================================

\subsection{Why Latitude and Longitude Break Area Calculations}
GPS coordinates use degrees (e.g., $22.58^\circ\,\text{N}, 88.41^\circ\,\text{E}$). 
\begin{itemize}
    \item Earth is a sphere, not a flat tabletop. 
    \item One degree of latitude is roughly $111\,\text{km}$, but one degree of longitude shrinks as you move away from the equator towards the poles.
    \item If you calculate tree area using raw degrees, your calculations will be completely distorted.
\end{itemize}

\subsection{Converting to UTM (Meters)}
To do real math, \texttt{Clover} converts GPS degrees into a local \textbf{UTM grid} (Universal Transverse Mercator, like \texttt{EPSG:32645} for Kolkata):
\begin{equation}
(X, Y) \text{ coordinates are strictly in meters.}
\end{equation}
Now, $1\,\text{unit} = 1\,\text{meter}$, and calculating area is as simple as length times width.

% ==============================================================================
\section{Step 3: The Big Image Problem (Sliding Window)}
% ==============================================================================

\subsection{Why Satellites Crash Computers}
A high-resolution satellite scene can be $10,000 \times 10,000$ pixels. In memory, that is over a gigabyte of raw numbers. If you feed that entire image into a neural network all at once, your computer immediately crashes with an \textbf{Out Of Memory (OOM)} error.

\subsection{The Puzzle-Piece Solution}
\texttt{Clover} cuts the huge image into small tiles (usually $400 \times 400$ pixels), like puzzle pieces:
\begin{itemize}
    \item \textbf{20\% Overlap Halo}: If a tree happens to sit right on the cut between two tiles, a naive slice would cut the tree in half and miss it.
    \item By overlapping tiles by 20\%, every single tree is guaranteed to land fully inside at least one tile.
    \item The AI runs on each small tile, and \texttt{Clover} maps the local pixel detections back to their real-world GPS/UTM coordinates on the master map.
\end{itemize}

% ==============================================================================
\section{Step 4: Finding the Trees (AI Model)}
% ==============================================================================

\subsection{DeepForest and RetinaNet}
To spot trees from above, \texttt{Clover} uses \textbf{DeepForest}, an open-source neural network built specifically for forestry:
\begin{itemize}
    \item \textbf{What it looks for}: The bright sunlit dome of the tree crown, the rounded circular texture of branches, and the dark shadow cast on the ground.
    \item \textbf{Pretrained on 30 Million Trees}: It has already seen aerial images of pine forests, tropical broadleaves, savannas, and urban parks.
\end{itemize}

\subsection{Focal Loss: How the AI Focuses}
In satellite imagery, $98\%$ of the image is just background: grass, dirt, roads, or roofs. Only $2\%$ is tree tops. 
\begin{itemize}
    \item If you train normal AI, it gets lazy: it simply guesses ``background'' everywhere and gets $98\%$ accuracy without learning anything!
    \item \textbf{Focal Loss} fixes this by automatically turning down the volume on easy background pixels and forcing the model to concentrate on the hard, ambiguous tree edges.
\end{itemize}

\subsection{Removing Duplicate Boxes (NMS)}
Because our tiles overlap by 20\%, the same tree near a tile boundary might get detected twice.
\begin{equation}
\text{IoU} = \frac{\text{Shared Area}}{\text{Total Combined Area}}
\end{equation}
If two detected boxes overlap with an $\text{IoU} > 15\%$, \texttt{Clover} throws away the lower-confidence duplicate and keeps the cleanest detection.

% ==============================================================================
\section{Step 5: The Secret Sauce: Canopy Geometry}
% ==============================================================================

This is where \texttt{Clover} does what other tools fail to do.

\subsection{The 21.5\% Bounding Box Error}
When the AI detects a tree, it outputs a rectangular box of width $W$ and height $H$:
\begin{equation}
\text{Box Area} = W \times H
\end{equation}
A real tree crown is an ellipse inscribed inside that box:
\begin{equation}
\text{Tree Area} = \pi \times \frac{W}{2} \times \frac{H}{2} = \frac{\pi}{4} \times W \times H \approx 0.7854 \times \text{Box Area}
\end{equation}

\begin{plainbox}{The Corner Penalty}
The four corners of a bounding box contain empty dirt, grass, or walkways. If you don't multiply by $\frac{\pi}{4}$ ($0.7854$), \textbf{you overestimate tree canopy by $21.46\%$ on every single tree}. \texttt{Clover} converts all boxes into true geometric ellipses.
\end{plainbox}

\subsection{Dissolving Overlapping Branches (Unary Union)}
In real life, trees in a park or forest grow close together and their branches interlock.
\begin{itemize}
    \item If Tree 1 covers $50\,\text{m}^2$ and Tree 2 covers $50\,\text{m}^2$, but they overlap by $20\,\text{m}^2$:
    \item \textbf{Naive addition}: $50 + 50 = 100\,\text{m}^2$ (Wrong! You counted the overlap twice).
    \item \textbf{Real canopy cover}: $50 + 50 - 20 = 80\,\text{m}^2$.
\end{itemize}

\texttt{Clover} uses a computational geometry operation called \textbf{Topological Unary Union} ($\bigcup \mathcal{P}_i$). Think of it like taking overlapping droplets of water: as soon as they touch, they melt into one single continuous puddle. This gives the exact, true green canopy cover over the ground.

\subsection{The Li et al. (2023) +20\% Calibration}
Satellites see the dense top leaves of a tree, but thin outer twigs and sunlit edge leaves are semi-translucent from orbit. 

A landmark study by \textbf{Li et al. (published in PNAS Nexus, 2023)} compared satellite AI detections against ground laser scans. They found that satellite models consistently underestimate crown area by about \textbf{20\%}. \texttt{Clover} applies this peer-reviewed $+20\%$ calibration factor so the final number matches ground reality.

% ==============================================================================
\section{Step 6: Turning Tree Size into Carbon (Allometry)}
% ==============================================================================

\subsection{How Do You Weigh a Tree Without Chopping It Down?}
You cannot put a live 30-meter tree on a weighing scale. Instead, foresters use \textbf{Allometry}: the mathematical relationship between the physical dimensions of a tree and its dry wood weight (Above-Ground Biomass, or AGB).

\subsection{The Jucker et al. (2016) Equation}
Forest ecologist Dr. Tommaso Jucker analyzed over 2,000 destructively weighed trees across the world and proved that crown diameter from above ($CD$) predicts wood weight:
\begin{equation}
\text{Biomass (kg)} = \exp\left( -0.328 + 2.404 \times \ln(CD) \right)
\end{equation}
For example:
\begin{itemize}
    \item A small $3\,\text{m}$ crown $\implies \approx 100\,\text{kg}$ of wood.
    \item A mature $7\,\text{m}$ crown $\implies \approx 780\,\text{kg}$ of wood.
    \item A massive $12\,\text{m}$ banyan $\implies \approx 3,500\,\text{kg}$ (3.5 tonnes) of wood.
\end{itemize}

\subsection{From Wood to Carbon and $\text{CO}_2$}
\begin{enumerate}
    \item \textbf{Dry wood to Carbon}: Wood is roughly $47\%$ carbon (cellulose and lignin). So:
    \begin{equation}
    \text{Carbon} = \text{Wood Weight} \times 0.47
    \end{equation}
    \item \textbf{Carbon to Sequestered $\text{CO}_2$}: Carbon dioxide ($\text{CO}_2$) is heavier than pure carbon because it has two oxygen atoms. Using the molecular weight ratio ($44 / 12$):
    \begin{equation}
    \text{CO}_{2\text{e}} = \text{Carbon} \times 3.664
    \end{equation}
\end{enumerate}

% ==============================================================================
\section{Step 7: Statistical Honesty (The Error Margin)}
% ==============================================================================

\subsection{The Single-Tree Problem}
If you try to guess the weight of \textbf{just one tree} from its crown, you could be wrong by $\pm 56.5\%$. Some trees are tall and skinny; others are hollow inside.

\subsection{The Central Limit Theorem (Why Forest Averages Work)}
In a stand of 80 trees, random individual variations balance each other out: one tree is lighter than expected, another is heavier. Under the \textbf{Central Limit Theorem}, random statistical error shrinks with the square root of the tree count:
\begin{equation}
\text{Statistical Error} = \frac{56.5\%}{\sqrt{N_{\text{trees}}}}
\end{equation}
For $N = 80$ trees, the random sampling error drops to just $\mathbf{\pm 6.3\%}$.

\subsection{The Systematic Error Floor}
However, the general allometric equation itself has a baseline regional error of about $\pm 12\%$ (because wood density varies across countries). You can never reduce your error below this floor.

\texttt{Clover} combines both errors honestly:
\begin{equation}
\text{Total Stand Error} = \sqrt{(\text{Statistical Error})^2 + (12\%)^2} \approx \mathbf{\pm 13.5\%}
\end{equation}
When \texttt{Clover} presents a report, it does not say ``This park has 5.55 tonnes of biomass.'' It says:
$$\mathbf{5.55 \text{ tonnes } \pm 0.75 \text{ tonnes (Range: } 4.80 \text{ to } 6.30 \text{ tonnes)}}$$
Carbon credit auditors love this because it is honest, reproducible, and verifiable.

% ==============================================================================
\section{Step 8: When Does the Tool Fail? (Honest Limitations)}
% ==============================================================================

Every physical tool has limits. \texttt{Clover} documents four specific failure modes:

\begin{enumerate}
    \item \textbf{Dense Closed Canopies}: In thick jungle (like core Sundarbans), trees touch shoulder-to-shoulder, creating an unbroken green ceiling. You cannot see individual tree boundaries. In this case, \texttt{Clover} uses continuous green canopy segmentation rather than pretending to count individual stems.
    \item \textbf{Understory Blindness}: Satellites only see the roof of the forest. If small bushes or young saplings are growing underneath large trees, an optical satellite cannot see them.
    \item \textbf{Shadows at Sunset/Sunrise}: When photos are taken early in the morning, trees cast long black shadows. Shadows can hide neighboring trees or be mistaken for dark leaves.
    \item \textbf{Winter / Leaf-Off}: Deciduous trees that shed their leaves in winter look like bare sticks from space, causing the AI to miss them.
\end{enumerate}

% ==============================================================================
\section{Quick Summary Table}
% ==============================================================================

\begin{table}[h]
\centering
\small
\begin{tabular}{@{}lll@{}}
\toprule
\textbf{Step} & \textbf{What We Do} & \textbf{Why We Do It} \\ \midrule
\textbf{1. GSD Check} & Validate resolution ($< 0.5\,\text{m}$) & Prevents running AI on blurry images \\
\textbf{2. UTM Projection} & Convert degrees to meters & Ensures real Euclidean square meters \\
\textbf{3. Sliding Window} & Slice into $400 \times 400$ with $20\%$ overlap & Prevents computer out-of-memory crashes \\
\textbf{4. DeepForest} & RetinaNet with Focal Loss & Spots trees and ignores empty grass/dirt \\
\textbf{5. Ellipse Shape} & Multiply box area by $\pi/4$ ($0.7854$) & Removes empty box corners ($-21.5\%$ error) \\
\textbf{6. Unary Union} & Dissolve overlapping polygons & Stops counting shared branches twice \\
\textbf{7. Li Calibration} & Add $+20\%$ to canopy area & Compensates for semi-translucent branch edges \\
\textbf{8. Allometry} & Jucker et al. power-law scaling & Converts crown width into wood weight \\
\textbf{9. CLT Bounds} & Propagate dual uncertainty & Provides an auditable $\pm 13\%$ error bracket \\ \bottomrule
\end{tabular}
\caption{The complete Clover pipeline in 9 simple steps.}
\end{table}

\end{document}
